\documentclass[10pt,a4paper]{amsart}
\usepackage[margin=19mm]{geometry}
\usepackage{amsmath,amssymb,mathtools}
\usepackage[scr=rsfs,cal=euler]{mathalfa}
\usepackage{xfrac}
\usepackage[unicode,hidelinks,bookmarks=false]{hyperref}
\newcommand{\ie}{i.e.\ }
\newcommand{\cf}{cf.\ }
\newcommand{\resp}{respectively\ }
\newcommand{\Subsection}{\subsection}
\providecommand{\nobreakdash}{\nobreak\hbox{-}}
\setlength{\emergencystretch}{2em}
\multlinegap0pt
\usepackage{bm}
\usepackage{bbm}
\usepackage{dsfont}
\usepackage{xspace}
\usepackage{cancel}
\usepackage{mathtools}
\usepackage{amsthm}
\makeatletter
\@ifundefined{thm}{\newtheorem{thm}{Thm}}{}
\@ifundefined{prop}{\newtheorem{prop}[thm]{Proposition}}{}
\makeatother
\makeatletter
\newcommand{\R}{{\mathbb R}}
\newcommand{\Z}{{\mathbb Z}}
\expandafter\ifx\csname proofmthm\endcsname\relax
\newenvironment{proofmthm}{\begin{proof}[Proof of the Main Theorem]}
                   {\end{proof}}
\fi
\expandafter\ifx\csname say\endcsname\relax
\newenvironment{say}{}{}
\fi
\expandafter\ifx\csname sketch\endcsname\relax
\newenvironment{sketch}{\begin{proof}[Sketch of proof]}{\end{proof}}
\fi
\makeatother
\title[Mld versus lct near zero]{Mld versus lct near zero}
\author{Florin Ambro}
\date{Source: arXiv:2506.18507v1 (CC BY 4.0)}
\begin{document}
\maketitle

\noindent\emph{Source and licence.} Mld versus lct near zero. Version arXiv:2506.18507v1 (CC BY 4.0), distributed under the Creative Commons Attribution 4.0 International licence, as recorded in the arXiv licence metadata for this exact version and shown on the abstract page https://arxiv.org/abs/2506.18507v1. Full rights notice: licenses/arxiv-2506.18507-CC-BY-4.0.txt.\par\smallskip

\noindent\textit{Editorial scope.} Counted unit: the Prop whose source label is \texttt{extpofcn} in the article (source0.tex lines 542--550, source label \texttt{extpofcn}), delivered together with the article's own complete original proof of it (source0.tex lines 552--630, one \texttt{proof} environment of 79 lines). No mathematics is written, repaired or re-derived: statement and proof are the article's own text line for line. Required context retained uncounted: none. Every definition, display and label of those lines is retained unchanged. Omitted from this package and not redistributed: 1--541, 551--551, 631--1430. Nothing inside a retained excerpt is omitted or altered: every retained body is delivered exactly as the article's lines are written, so no omission receipt of any kind accompanies this package. Citations: the retained excerpts carry no citation command at all, so no bibliography entry of the article is transcribed and no reference is redistributed. Reference adaptation: none was needed and none was made; every reference inside the delivered unit resolves inside it, so no unresolved reference or citation is produced. Printed slips kept literal: no printed word, symbol, hypothesis or number of the counted statement or of its proof was altered. Layout and class: the article's own class is \texttt{amsart} with packages including amsthm, amsfonts, amssymb, amscd, euscript, xy; the wrapper is the lane's pinned \texttt{amsart} editorial wrapper at 10pt on A4, which restates the theorem-like environments the retained text needs and adds the article's own macro definitions that the retained text needs (\texttt{\textbackslash R}, \texttt{\textbackslash Z}), with extra wrapper packages (bm, bbm, dsfont, xspace, cancel, mathtools). The delivered numbering is the unit's own. Assets: no retained span contains a figure environment, a graphics-inclusion command, a PDF-page inclusion, a table or a TikZ picture; no image file is copied into this package, the package declares no asset dependency and no image file is redistributed with it. Licence: version arXiv:2506.18507v1 of this article is distributed under the Creative Commons Attribution 4.0 International licence, as recorded in the arXiv licence metadata for this exact version and shown on the abstract page https://arxiv.org/abs/2506.18507v1. A full rights notice is delivered at licenses/arxiv-2506.18507-CC-BY-4.0.txt. Characters: every retained excerpt is pure ASCII, so the delivered engine, XeLaTeX with its default Latin Modern text font, reports no missing character.
\begin{prop}\label{extpofcn}
	Let $\varphi_0\colon \Lambda_0\to \Z$ be a lattice homomorphism such that
	$\varphi_0(C_0)=[0,l_0]$ for some $l_0>0$. Then there exists a lattice
	homomorphism $\varphi'\colon \Lambda\to \Z$ such that:
	\begin{itemize}
		\item $\varphi'|_{\Lambda_0}=q\varphi_0$ for some integer $1\le q< w$.
		\item $\varphi'(C)=[0,l']$ for some $0<l'\le wl_0$.
	\end{itemize}
\end{prop}

\begin{proof} The inclusion of lattices $j\colon \Lambda_0\subset \Lambda$ is a direct
	summand. It dualizes to a lattice projection $j^*\colon \Lambda^*\to \Lambda_0^*$,
	which identifies $\Lambda_0^*$ with the quotient $\Lambda^*/\Z \varphi$. 
	Let $C^*\subset \Lambda^*\otimes_\Z \R$ be the convex set polar dual to $C\subset \Lambda_\R$. 
	Let $C^*_0\subset \Lambda_0^*\otimes_\Z \R$ be the $j^*$-image of 
	$C\subset \Lambda^*\otimes_\Z \R$. Then $C_0\subset \Lambda_0\otimes_\Z \R$ is polar 
	dual to $C_0^*\subset \Lambda^*_0\otimes_\Z \R$. 
	
	Let $\varphi(C)=[-w_-,w_+]$, with $w_-,w_+>0$ and $w_-+w_+=w$.
	Denote $\varphi_1=\varphi$. Let $\varphi_2\colon \Lambda\to \Z$ be an arbitrary extension of 
	$\varphi_0$. Then $\varphi_1,\varphi_2\in \Lambda^*$ are linearly independent in $\Lambda^*_\R$.
	
	 By assumption, $\varphi_0$ is non-negative 
    on $\cup_{t>0}tC_0=\sigma\cap \varphi^\perp$. Equivalently, there exists $c\in \R$ such that 
	$$
	c\varphi_1+\varphi_2\in \sigma^\vee.
	$$
	By assumption, there exist finitely many elements $e_i \in \Lambda^{prim}\cap C$ such that
	$\sigma=\sum_i\R_{\ge 0}e_i\subset \Lambda_\R$. The condition for $c\varphi_1+\varphi_2\in \sigma^\vee$ 
	becomes
	$$
	\max_{\varphi_1(e_i)>0}\frac{-\varphi_2(e_i)}{\varphi_1(e_i)}\le c\le
	\min_{\varphi_1(e_i)<0}\frac{\varphi_2(e_i)}{-\varphi_1(e_i)}.
	$$
	Denote $\gamma'=\frac{1}{l_0}$. 
	Since $-\gamma'\varphi_0 \in C^*_0$, there exists $t\in \R$ such that 
	$t\varphi_1-\gamma'\varphi_2\in C^*$. Denote $c'=\gamma'c+t$. Then 
	$$
	c'\varphi_1-\gamma'(c\varphi_1+\varphi_2)\in C^*.
	$$
	
	{\em Step 1}: We claim that $-\frac{1}{w_+}\le c'\le \frac{1}{w_-}$. Indeed, from $\sigma^\vee\subseteq C^*$
	we deduce
	 $$
	c'\varphi_1=
	(c'\varphi_1-\gamma'(c\varphi_1+\varphi_2))+\gamma'(c\varphi_1+\varphi_2)
	\in C^*.
	$$
	The assumption $\varphi(C)=[-w_-,w_+]$ is polar dual to 
	$\{x\in \R | x \varphi_1\in C^* \}=[-\frac{1}{w_+},\frac{1}{w_-}]$. Therefore the claim holds.
	
	{\em Step 2}: We claim that $-\gamma'\frac{\min(w_-,w_+)}{w_-+w_+}\cdot (c\varphi_1+\varphi_2)\in C^*$.
	
	Indeed, suppose $c'\le 0$. Then 
	$$
	-\epsilon\gamma'(c\varphi_1+\varphi_2)=\epsilon(c'\varphi_1-\gamma'(c\varphi_1+\varphi_2))+(1-\epsilon)\frac{\varphi_1}{w_-}\in C^*,
	$$
	where $\epsilon=\frac{1}{1-c'w_-}\in (0,1]$. Since $c'\ge -\frac{1}{w_+}$, we deduce $\epsilon\ge  \frac{w_+}{w_-+w_+}$.
	
	Suppose $c'\ge 0$. Then 
	$$
	-\epsilon\gamma'(c\varphi_1+\varphi_2)=\epsilon(c'\varphi_1-\gamma'(c\varphi_1+\varphi_2))+(1-\epsilon)\frac{-\varphi_1}{w_-}\in C^*,
	$$
	where $\epsilon=\frac{1}{1+c'w_+}\in (0,1]$.
	Since $c'\le \frac{1}{w_-}$, we deduce $\epsilon\ge  \frac{w_-}{w_-+w_+}$.
	
	{\em Step 3}:
	Suppose $w_-\le w_+$. Choose $c=\min_{\varphi_1(e_i)<0}\frac{\varphi_2(e_i)}{-\varphi_1(e_i)}$.
	From Step 2 we have
	$$
	-\gamma'\frac{w_-}{w_-+w_+}\cdot (c\varphi_1+\varphi_2)\in C^*.
	$$
	There exists $e_i$ such that $\varphi_1(e_i)<0$ and $c=\frac{\varphi_2(e_i)}{-\varphi_1(e_i)}$.
	Here $\varphi_1(e_i)=-q$ for some positive integer $q$. Set 
	$\varphi'=\varphi_2(e_i)\varphi_1-\varphi_1(e_i) \varphi_2 \in \Lambda\setminus 0$.
	Then $\varphi'=q(c\varphi_1+\varphi_2)\in \sigma^\vee$, $j^*(\varphi')=q\varphi_0$ and
	$$
	-\gamma'\frac{w_-}{w_-+w_+}\frac{1}{q}\varphi' \in C^*.
	$$
	Since $e_i \in C$ and $\varphi_1(C)\subset [-w_-,w_+]$, 
	we obtain $q\le w_-$. Therefore $q<w$ and 
	$$
	\frac{w_-}{w_-+w_+}\frac{1}{q} \ge  \frac{1}{w}.
	$$
	We obtain $-\frac{\gamma'}{w}\varphi'\in C^*$. Then $\varphi'(C)$ is contained
	in $[0,\frac{w}{\gamma'}]=[0,wl_0]$.
	
	If $w_+\le w_-$ we choose $c$ to be the other extremal point, and argue similarly.
\end{proof}

\end{document}
